\documentclass[12pt,compress,aspectratio=169]{beamer}
\input{../../mybeamer}

\usetikzlibrary{decorations.pathmorphing,patterns}


\title{Class 5: Momentum, Impulse and Collisions}
\subtitle{AP Physics C--Mechanics}
\input{../../me}
\input{../../mycommands}

\begin{document}

\begin{frame}
  \maketitle
\end{frame}



\section{Momentum}

\begin{frame}{Momentum}
  \textbf{Momentum} (or \textbf{translational momentum}, or \textbf{linear
    momentum}) quantifies motion: \emph{how much} (mass) is moving; and
  \emph{how fast} (velocity) it moves.

  \eq{-.1in}{
    \boxed{\vec p=m\vec v}
  }
  \begin{center}
    \begin{tabular}{l|c|c}
      \rowcolor{pink}
      \textbf{Quantity}   & \textbf{Symbol} & \textbf{SI Unit} \\ \hline
      Momentum  & $\vec p$ & \si{\kilo\gram\metre\per\second} \\
      Mass      & $m$     & \si{\kilo\gram} \\
      Velocity  & $\vec v$ & \si{\metre\per\second}
    \end{tabular}
  \end{center}
  For rotational motion of rigid bodies, there is also \textbf{angular momentum}
  which will be studied in in Classes \#7, \#9 and \#10.
\end{frame}



\begin{frame}{General Form of First \& Second Laws of Motion}
  Taking the time derivative of the momentum vector from an inertial frame of
  reference using the chain rule:

  \eq{-.1in}{
    \diff{\vec p}t=\diff{(m\vec v)}t
    =m\diff{\vec v}t+\diff mt\vec v
    =m\vec a+\dot m\vec v
  }

  For and object with constant mass $m$ (i.e.\ $\dot m=0$), the right-hand side
  of the equation reduces to the familiar form of the second law of motion,
  $m\vec a$.
\end{frame}




\begin{frame}{General Form of First \& Second Laws of Motion}
  Accounting for changing mass, the general form of the first and second law:
  \begin{enumerate}[I.]
  \item\textbf{The momentum of an object or a system of objects remains
    constant until a net external force acts on it.}
  \item\textbf{The net external force acting on an object or a system of
    objects is the rate of change of its momentum.}
  \end{enumerate}
  Summarizing the first two laws into a single equation:

  \eq{-.1in}{
    \boxed{
      \vec F_\text{net}(t)=\diff{\vec p}t
    }
  }

  %Like the ``special case'' from the dynamics section,
  This equation is only applicable from an inertial frame of reference.
\end{frame}



\section{Impulse}

\begin{frame}{Impulse}
  \textbf{Impulse} $\vec J$ is defined as the time integral of force $\vec F$:

  \eq{-.1in}{
    \vec J=\int_{t_0}^{t_1}\vec F(t)\dl t
  }

  We can calculate the impulse generated by
  \begin{itemize}
  \item any of the forces acting on the object, or
  \item the net force (called the \textbf{net impulse})
  \end{itemize}
  over the time interval between $t_0$ and $t_1$. Since $\vec F$, $\vec p$ and
  $\vec J$ are all vectors, the integral can be evaluated independently in each
  of the $\iii$, $\jjj$ and $\kkk$ directions:
  
  \eq{-.1in}{
    J_x=\int_{t_0}^{t_1}F_x(t)\dl t\quad\quad
    J_y=\int_{t_0}^{t_1}F_y(t)\dl t\quad\quad
    J_z=\int_{t_0}^{t_1}F_z(t)\dl t
  }
\end{frame}



\begin{frame}{Impulse-Momentum Theorem}
  Rearranging the variables in the general form of the second law of motion:
  
  \eq{-.1in}{
    \vec F_\text{net}(t)=\diff{\vec p}t\quad\to\quad
    \vec F_\text{net}(t)\dl t=\dl\vec p
  }

  Integrating both sides, we get the \textbf{impulse-momentum theorem}:
 
  \eq{-.1in}{
    \int_{t_0}^{t_1}\vec F_\text{net}(t)\dl t
    =\int_{\vec p_0}^{\vec p_1}\dl\vec p
    \quad\to\quad
    \boxed{ \vec J_\text{net} =\Delta\vec p}
  }
  
  where $\vec p_0=\vec p(t_0)$ and $\vec p_1=\vec p(t_1)$. The
  \textbf{net impulse} is equal to the change in momentum of the object or
  system of objects.
\end{frame}



\begin{frame}{Average Force}
  \textbf{Average force} $\vec F_\text{avg}$ is the time-averaged force vector
  that generated the same impulse. It is used extensively in introductory
  physics courses to avoid integration:

  \eq{-.1in}{
    \vec J=\int_{t_0}^{t_1}\vec F(t)\dl t=\vec F_\text{avg} \Delta t
  }

  \vspace{.2in}\textbf{Note:} The average force above is a \emph{time}
  averaging, like we have derived average velocity and average acceleration in
  kinematic in Class \#1. This is different from the \emph{spatial} averaging
  for calculating work in Classes \#3 and \#4.
\end{frame}



\begin{frame}{Impulse: An Example}
  \textbf{Example:} Jim pushes a box with mass \SI{1.0}{\kilo\gram} with a
  \SI{5.0}{\newton} force for \SI{10}{\second} while the box stays on the same
  place. Find the impulse of the pushing force, friction force, the
  gravitational force, and the net force.
\end{frame}



\begin{frame}{Rocket Propulsion Problem}
  \textbf{Example:} A rocket generates a thrust force by ejecting hot gases
  from an engine. If it takes \SI1{\milli\second} to combust
  \SI{1.0}{\kilo\gram} of fuel, ejecting it at a speed of
  \SI{1000}{\metre\per\second}, what thrust is generated?
  \begin{enumerate}[(A)]
  \item \SI{1000}\newton
  \item \SI{10000}\newton
  \item \SI{100000}\newton
  \item \SI{1000000}\newton
  \end{enumerate}
\end{frame}



\begin{frame}{Another Space Example}
  \textbf{Example:} A rocket for mining the asteroid belt is designed like a
  large scoop. It is approaching asteroids at a velocity of
  \SI{e4}{\metre\per\second}. The asteroids are much smaller than the rocket.
  If the rocket scoops asteroids at a rate of \SI{100}{\kilo\gram\per\second},
  what thrust (force) must the rocket's engine provide in order for the rocket
  to maintain constant velocity? Ignore any variation in the rocket's mass due
  to the burning fuel.
  \begin{enumerate}[(A)]
  \item\SI{e3}\newton
  \item\SI{e6}\newton
  \item\SI{e9}\newton
  \item\SI{e12}\newton
  \end{enumerate}
\end{frame}


%\begin{frame}{Impulse: Another Example}
%  \textbf{Example:} Two balls of the same mass are dropped from the same
%  height onto the floor. The first ball bounces upwards from the floor
%  elastically. The second ball sticks to the floor. The first applies an
%  impulse to the floor of $I_1$ and the second applies an impulse $I_2$. The
%  two impulses obey:
%  \begin{enumerate}[(a)]
%  \item $I_2=2I_1$
%  \item $I_2=I_1/2$
%  \item $I_2=4I_1$
%  \item $I_2=I_1/4$
%  \end{enumerate}
%\end{frame}


\begin{frame}{Momentum and Kinetic Energy}
  We can express kinetic energy ($K$) in terms of the magnitude of momentum
  ($p=|\vec p|$):
  
  \eq{-.1in}{
    K
    =\frac12mv^2
    =\frac12\left[\frac{\color{magenta}m^2}m\right]{\color{magenta}v^2}
    \quad\to\quad
    \boxed{
      K=\frac{\color{magenta}p^2}{2m}
    }
  }
  \begin{itemize}
  \item Whenever momentum changes in \emph{magnitude}, kinetic energy also
    changes
  \item\textbf{BUT:} Kinetic energy $K$ is a scalar quantity, but momentum
    $\vec p$ is the vector, therefore
  \item A momentum vector that changes only in \emph{direction} does not
    change in kinetic energy, which means that
  \item A net force may impart a net impulse ($\vec J_\text{net}$) on the
    object, while not doing any net work ($W_\text{net}$) on it
  \end{itemize}
\end{frame}



\begin{frame}{Before We Dive Into Some Exercises}
  Typical applications of momentum conservation are collision and explosion
  problems
  \begin{itemize}
  \item\textbf{Collision: A hits B}
    \begin{itemize}
    \item Regardless of whether they move together or not afterwards, momentum
      is conserved
    \item Head-on collisions are usually 1D
    \item Glancing collisions are usually 2D or 3D
    \end{itemize}
  \item\textbf{Explosion: A explodes and becomes B and C (and D and E\ldots)}
    \begin{itemize}
    \item A perfectly inelastic collision in reverse
    \item Total momentum of B and C (and D and E\ldots) is the same as A in the
      beginning
    \item Usually a 2D or 3D problem
    \end{itemize}
  \end{itemize}
\end{frame}



\section{Collisions}

\begin{frame}{Conservation of Momentum}
  When two objects collide in isolation\footnote{i.e.\ The collision is an
  isolated system. We have studied this concept in dynamics (Class \#2) and
  also in conservation of energy (Classes \#3 and \#4). During the collision,
  there are no external forces}, from the third law of motion, we know that the
  action and reaction forces between the two objects are equal in magnitude and
  opposite in direction.
  \begin{center}
  \begin{tikzpicture}[scale=1.1]
    \draw[thick,fill=lightgray] circle (.3);
    \draw[axes,gray] (-1.5,0)--(-.4,0);
    \draw[gray,fill=gray!20,dash dot] (-1.8,0) circle (.3) node{$m_1$};
    
    \begin{scope}[rotate=-35]
      \draw[thick,fill=lightgray] (.58,0) circle (.3);
      \draw[axes,gray] (2,0)--(1.4,0);
      \draw[gray,fill=gray!20,dash dot] (2.3,0) circle (.3) node{$m_2$};
      \draw[ultra thick] (.29,.1)--(.29,-.1);
      \draw[vectors,magenta] (.29,0)--+(1,0) node[pos=1.2]{$F_n$};
      \draw[vectors,red] (.29,0)--+(-1,0) node[pos=1.2]{$F_n$};
      \draw[vectors,blue](.27,0)--+(0,.7) node[pos=1.2]{$F_f$};
      \draw[vectors,cyan](.31,0)--+(0,-.7) node[pos=1.2]{$F_f$};
    \end{scope}
    %\fill[red] circle (.03);
  \end{tikzpicture}
\end{center}

  The forces the objects apply to each other may include (but not limited to):
  normal force ($F_n$), static or kinetic friction ($\vec F_f$), elastic force
  ($F_e$), electrostatic force ($F_q$) if they are charged, or magnetic force
  ($F_m$) if they are magnets.
  \vspace{.2in}
\end{frame}



\begin{frame}{Conservation of Momentum}
  \begin{center}
  \begin{tikzpicture}[scale=1.1]
    \draw[thick,fill=lightgray] circle (.3);
    \draw[axes,gray] (-1.5,0)--(-.4,0);
    \draw[gray,fill=gray!20,dash dot] (-1.8,0) circle (.3) node{$m_1$};
    
    \begin{scope}[rotate=-35]
      \draw[thick,fill=lightgray] (.58,0) circle (.3);
      \draw[axes,gray] (2,0)--(1.4,0);
      \draw[gray,fill=gray!20,dash dot] (2.3,0) circle (.3) node{$m_2$};
      \draw[ultra thick] (.29,.1)--(.29,-.1);
      \draw[vectors,magenta] (.29,0)--+(1,0) node[pos=1.2]{$F_n$};
      \draw[vectors,red] (.29,0)--+(-1,0) node[pos=1.2]{$F_n$};
      \draw[vectors,blue](.27,0)--+(0,.7) node[pos=1.2]{$F_f$};
      \draw[vectors,cyan](.31,0)--+(0,-.7) node[pos=1.2]{$F_f$};
    \end{scope}
    %\fill[red] circle (.03);
  \end{tikzpicture}
\end{center}

  Since the net force on $m_1$ and $m_2$ are equal and opposite, and both
  objects experience the net force over the same time interval, the net
  impulse on $m_1$ and $m_2$ are also equal and opposite, and the change in
  momentum on $m_1$ and $m_2$ are also equal and opposite:

  \eq{-.1in}{
    \vec F_1=-\vec F_2\quad\to\quad
    \vec J_1=-\vec J_2\quad\to\quad
    \boxed{
      \Delta\vec p_1=-\Delta\vec p_2
    }
  }
\end{frame}



\begin{frame}{Conservation of Momentum}
  The total \emph{change} in momentum is therefore zero:

  \eq{-.1in}{
    \boxed{
      \Delta\vec p_1+\Delta\vec p_2=0
    }
  }
    
  In other words, during a collision, the \textbf{total momentum of the system
  is conserved during a collision}:
    
  \eq{-.1in}{
    \boxed{
      \sum_i\vec p_i=\sum_i\vec p_i'
    }
  }

  The total momentum before the collision is the same as the total momentum
  after the collision
\end{frame}




\begin{frame}{Conservation of Momentum}
  \begin{center}
  \begin{tikzpicture}[scale=1.1]
    \draw[thick,fill=lightgray] circle (.3);
    \draw[axes,gray] (-1.5,0)--(-.4,0);
    \draw[gray,fill=gray!20,dash dot] (-1.8,0) circle (.3) node{$m_1$};
    
    \begin{scope}[rotate=-35]
      \draw[thick,fill=lightgray] (.58,0) circle (.3);
      \draw[axes,gray] (2,0)--(1.4,0);
      \draw[gray,fill=gray!20,dash dot] (2.3,0) circle (.3) node{$m_2$};
      \draw[ultra thick] (.29,.1)--(.29,-.1);
      \draw[vectors,magenta] (.29,0)--+(1,0) node[pos=1.2]{$F_n$};
      \draw[vectors,red] (.29,0)--+(-1,0) node[pos=1.2]{$F_n$};
      \draw[vectors,blue](.27,0)--+(0,.7) node[pos=1.2]{$F_f$};
      \draw[vectors,cyan](.31,0)--+(0,-.7) node[pos=1.2]{$F_f$};
    \end{scope}
    %\fill[red] circle (.03);
  \end{tikzpicture}
\end{center}

  Of course, this is to be expected. From the \emph{first} law of motion, when
  there is no net \emph{external} force in a system, the total momentum of the
  system is constant.
  \begin{itemize}
  \item The forces that $m_1$ and $m_2$ apply on each other are \emph{internal}
    forces
  \item The third law of motion must be true because it is an application of
    the first law of motion
  \end{itemize}
\end{frame}


\begin{frame}{Conservation of Momentum}
  For a collision of two objects\footnote{Or an explosion}, momentum is
  conserved at the ime of interaction:
  
  \eq{-.1in}{
    \boxed{\vec p_1+\vec p_2=\vec p_1'+\vec p_2'}
    \quad\rightarrow\quad
    \boxed{m_1\vec v_1+m_2\vec v_2=m_1\vec v_1'+m_2\vec v_2'}
  }
  \begin{center}
    \begin{tabular}{l|c|c}
      \rowcolor{pink}
      \textbf{Quantity} & \textbf{Symbol} & \textbf{SI Unit} \\ \hline
      Masses of 2 and 1 & $m_1$, $m_2$ & \si{\kilo\gram} \\
      Initial velocities of 1 and 2 & $\vec v_1$, $\vec v_2$ &
      \si{\metre\per\second} \\
      Final velocities of 1 and 2 & $\vec v_1'$, $\vec v_2'$ &
      \si{\metre\per\second}
    \end{tabular}
  \end{center}
  \textbf{Important note:} since momentum is a vector, the collision/explosion
  problem \emph{must} be solve using vector arithmetic.
\end{frame}



\begin{frame}{Glancing Collision in 2D}
  \vspace{.3in}\textbf{Example:} A billiard ball of mass \SI{.155}{\kilo\gram}
  (``cue ball'') moves with a velocity of \SI{1.25}{\metre\per\second} towards
  a stationary billiard ball (``eight ball'') of identical mass and strikes it
  with a glancing blow. The cue ball moves off at an angle of \ang{29.7}
  clockwise from its original direction, with a speed of
  \SI{.956}{\metre\per\second}. What is the final velocity of the eight ball?
\end{frame}



\section{Elastic Collisions}

\begin{frame}{Classifications of Collisions}
  \begin{itemize}
  \item Elastic Collision:
    \begin{itemize}
    \item Total kinetic energy is conserved
    \item<alert@1> Momentum is conserved
    \end{itemize}
  \item Inelastic collision:
    \begin{itemize}
    \item Kinetic energy is \textbf{not} conserved
    \item<alert@1> Momentum is conserved
    \end{itemize}
  \item Completely inelastic collision:
    \begin{itemize}
    \item ``Perfectly inelastic collision''
    \item A special case of inelastic collision
    \item The objects move together after the collision
    \item Kinetic energy is \textbf{not} conserved
    \item<alert@1> Momentum is conserved
    \end{itemize}
  \end{itemize}
\end{frame}



\begin{frame}{Elastic Collisions}
  \begin{center}
  \begin{tikzpicture}[scale=1.1]
    \draw[thick,fill=lightgray] circle (.3);
    \draw[axes,gray] (-1.5,0)--(-.4,0);
    \draw[gray,fill=gray!20,dash dot] (-1.8,0) circle (.3) node{$m_1$};
    
    \begin{scope}[rotate=-35]
      \draw[thick,fill=lightgray] (.58,0) circle (.3);
      \draw[axes,gray] (2,0)--(1.4,0);
      \draw[gray,fill=gray!20,dash dot] (2.3,0) circle (.3) node{$m_2$};
      \draw[ultra thick] (.29,.1)--(.29,-.1);
      \draw[vectors,magenta] (.29,0)--+(1,0) node[pos=1.2]{$F_n$};
      \draw[vectors,red] (.29,0)--+(-1,0) node[pos=1.2]{$F_n$};
      \draw[vectors,blue](.27,0)--+(0,.7) node[pos=1.2]{$F_f$};
      \draw[vectors,cyan](.31,0)--+(0,-.7) node[pos=1.2]{$F_f$};
    \end{scope}
    %\fill[red] circle (.03);
  \end{tikzpicture}
\end{center}

  When two objects collide in isolation (i.e.\ an isolated system with no
  external work $W_\text{net}=0$), total system energy is conserved. Within
  this system, work done by
  \begin{itemize}
  \item Conservative forces transform kinetic energies into potential energies
    (and from potential to kinetic energies)
  \item Non-conservative forces transform kinetic energies into thermal
    energies, and/or other forms of kinetic energies
  \end{itemize}
\end{frame}



\begin{frame}{Elastic Collisions}
  If the forces that two objects apply to each other are entirely
  conservative\footnote{Or if there are non-conservative forces, but those
  forces do not do any work during the collision}, then
  \begin{itemize}
  \item At the start of the collision, when two objects come together,
    conservative forces do negative work: kinetic energy is transformed into
    potential energies
  \item The the end of the collision, when the two objects move away from each
    other, conservative forces do positive work: potential energies are
    transformed back into kinetic energy
  \end{itemize}
  Therefore, during an elastic collision, kinetic energy is conserved during
  the collision, in addition to the conservation of momentum:

  \eq{-.1in}{
    \boxed{
      \sum_i K_i=\sum_i K_i'
    }\quad\text{\normalsize and}\quad
    \boxed{
      \sum_i \vec p_i=\sum_i \vec p_i'
    }
  }
  \vspace{.15in}
\end{frame}



%\begin{frame}{How to Solve Conservation of Momentum Problem}
%  \begin{enumerate}
%  \item Check whether the condition for the conservation of momentum is
%    satisfied (i.e.\ are there any external forces?)
%  \item If so, write out expressions for initial momentum and final momentum,
%    and equate the two. You will get 1 to 3 equations (one for each direction).
%  \item Solve these equations, find the quantity you need to find.
%  \end{enumerate}
%  Remember that momentum is a vector. If there is no external force component
%  in some direction, then the momentum component in this direction is still
%  conserved.
%\end{frame}



\begin{frame}{Conservative vs.\ Non-Conservative Forces}
  Bear in mind the list of conservative forces is \emph{very} short:
  \begin{itemize}
  \item Gravitational force
  \item Elastic force (only applied to \emph{ideal} springs)
  \item Electromagnetic forces
  \item Nuclear forces (only understood in quantum mechanics, not part of AP
    Physics C)
  \end{itemize}
  This means that most collisions must be inelastic
\end{frame}



\begin{frame}{Elastic Collisions}
  When the masses collide, kinetic energy is transformed into
  elastic potential energy:
  \begin{center}
    \begin{tikzpicture}
      \draw[thick,fill=magenta!30] (-1.8,0) circle (.3) node{$m_1$};
      \draw[thick,
        decoration={aspect=.5,segment length=3, amplitude=3, coil},
        decorate] (-1.5,0)--+(.8,0);
      \draw[vectors] (-.7,0)--+(1,0) node[midway,above]{$\vec v_1$};
      
      \draw[mass] (2.3,0) circle (.3) node{$m_2$};
      \draw[vectors] (2,0)--+(-1,0) node[midway,above]{$\vec v_2$};
    \end{tikzpicture}
  \end{center}

  When the charged particles collide, kinetic energy is
  transformed into electrical potential energy:
  \begin{center}
    \begin{tikzpicture}
      \draw[thick,fill=magenta!30] (-1.8,0) circle (.3) node{$q_1$};
      \draw[vectors] (-1.5,0)--+(1,0) node[midway,above]{$\vec v_1$};
      
      \draw[mass] (2.3,0) circle (.3) node{$q_2$};
      \draw[vectors] (2,0)--+(-1,0) node[midway,above]{$\vec v_2$};
    \end{tikzpicture}
  \end{center}

  When the masses ``collide'', kinetic energy is transformed
  into gravitational potential energy:
  \begin{center}
    \begin{tikzpicture}[thick,scale=1.2]
      \draw (1,0)--+(6,0);
      \draw[fill=magenta!30] (2,0)--(2,.5)--(2.5,.5) to[out=270,in=10] (2,0);
      \draw[vectors] (2.43,.25)--+(1,0) node[midway,above]{$\vec v_1$};
      
      \draw[mass] (4,0) to[out=10,in=270] (5.5,1)--(6,1)--(6,0)--(4,0);
      \node at (5.7,.5) {$m_2$};
      \node at (2.2,.3) {$m_1$};
      \draw[vectors] (5.05,.3)--+(-1,0) node[midway,above]{$\vec v_2$};
    \end{tikzpicture}
  \end{center}  
\end{frame}


%\begin{frame}{Example}
%  \textbf{Example:} Two blocks A and B, both have mass \SI{1.}{\kilo\gram}.
%  Block A has velocity \SI{3.}{\metre\per\second} and block B is at rest. Their
%  distance is \SI{1.}{\metre}. The surface is has dynamic friction coefficient
%  $\mu_k=0.02$. After they collide, they move together, what would be the final
%  velocity of these two blocks? How far can they go after the collision?
%\end{frame}



%\begin{frame}{More Example}
%  Max throws a ball into the air with an initial speed $\SI{10}{\metre\per\second} at an
%  angle of $60$ degree with the horizontal direction. By accident, the ball
%  splits into two parts (horizontally) in the air. Suppose both parts land at
%  the same time, neglecting the air resistance,
%  \begin{enumerate}
%  \item If one part is \SI{5}{\m} away from its original position (same
%    direction as the initial speed), where is the second part?
%  \item How about one parties \SI{5}{\metre} away from the original position in the
%    direction that has an angle of $30$ degree with its initial speed?
%  \end{enumerate}
%\end{frame}







%\begin{frame}{Collision Problem}
%  \textbf{Example:} Two objects with equal mass are heading towards each other
%  with equal speeds, undergo a head-on collision. Which one of the following
%  statement must be correct?
%  \begin{enumerate}[(A)]
%  \item Their final velocities are zero
%  \item Their final velocities may be zero
%  \item Each must have a final velocity equal to the other's initial velocity
%  \item Their velocities must be reduced in magnitude
%  \end{enumerate}
%\end{frame}
%
%
%
%\begin{frame}{Conservation of Momentum Example}
%  \textbf{Example:} Two astronauts, each of mass \SI{75}{\kilo\gram}, are
%  floating next to each other in space, outside the space shuttle. One of them
%  pushes the other through a distance of \SI{1.}{\metre} (about an arm's
%  length) with a force of \SI{300}{\newton}. What is the final relative
%  velocity of the two?
%
%  \vspace{.15in}\begin{enumerate}[A.]
%  \item \SI{2.}{\metre\per\second}
%  \item \SI{2.83}{\metre\per\second}
%  \item \SI{4.}{\metre\per\second}
%  \item \SI{16.}{\metre\per\second}
%  \end{enumerate}
%\end{frame}



%\begin{frame}{Example}
%  \textbf{Example:} A ball is dropped from a height $h$. It hits the ground
%  and bounces up with a momentum loss of \SI{10}{\percent} due to the impact.
%  The maximum height it will reach is:
%  \begin{enumerate}[(a)]
%  \item\num{.90}$h$
%  \item\num{.81}$h$
%  \item\num{.949}$h$
%  \item\num{.3}$h$
%  \end{enumerate}
%\end{frame}
%
%
%
%\begin{frame}{Conservation of Energy Example}
%  \textbf{Example:} A simple pendulum has a bob of mass \SI{2}{\kilo\gram}
%  hanging on a cord of length \SI{1}{\metre}. Suppose the pendulum is raised
%  until it is horizontal (and angular displacement of \ang{90}) and then
%  released. What is the speed of the bob at the bottom of its swing?
%  \begin{enumerate}[(a)]
%  \item\SI{9.91}{\metre\per\second}
%  \item\SI{19.6}{\metre\per\second}
%  \item\SI{3.13}{\metre\per\second}
%  \item\SI{4.43}{\metre\per\second}
%  \end{enumerate}
%\end{frame}
%
%
%
%\begin{frame}{Conservation of Energy Example}
%  \textbf{Example:} A toy firing a ball vertically consists of a vertical
%  spring which is compressed by \SI{.10}\metre. A force of \SI{10.0}{\newton}
%  is needed to hold the spring at that compression. If a ball of mass
%  \SI{.050}{\kilo\gram} is placed on the compressed spring and the spring is
%  released, the ball will reach a height (above its initial position) of:
%  \begin{enumerate}[(a)]
%  \item \SI{1.}{\metre}
%  \item \SI{1.2}{\metre}
%  \item \SI{1.4}{\metre}
%  \item \SI{1.6}{\metre}
%  \end{enumerate}
%\end{frame}







\begin{frame}{Derivation of Elastic Collision Equations}
  In elastic collisions, \emph{both} momentum and kinetic energy are conserved.
  In a 1D collision of two objects, we can rewrite the conservation of momentum
  by moving the $m_1$ terms to the left, and the $m_2$ to the right:

  \vspace{-.3in}{\large
    \begin{align}
      \nonumber
      m_1v_1 + m_2v_2&=m_1v_1'+m_2v_2'\\
      m_1(v_1-v_1')  &= m_2(v_2'-v_2)
    \end{align}
  }

  While the kinetic energy equation can be re-written in a similar form:
  
  \vspace{-.3in}{\large
    \begin{align}
      \nonumber
      \frac12 m_1v_1^2 + \frac12 m_2v_2^2 &=\frac12 m_1v_1'^2+\frac12 m_2v_2'^2\\
      m_1(v_1^2 -v_1'^2) &=m_2(v_2'^2-v_2^2)
    \end{align}
  }
\end{frame}



%\begin{frame}{Conservation of Momentum \& Energy in Elastic Collisions}
%  For collision of two objects, conservation of momentum equation can be
%  expressed as:
%
%  \vspace{-.1in}{\large
%    \begin{equation}
%      \boxed{m_1(v_1-v_1')=m_2(v_2'-v_2)}
%    \end{equation}
%  }
%
%  By moving $m_1$ terms to the left, and $m_2$ terms to the right. Likewise,
%  the conservation of energy can also be arranged as:
%
%  \vspace{-.1in}{\large
%    \begin{equation}
%      \boxed{m_1(v_1^2-v_1'^2)=m_2(v_2'^2-v_2^2)}
%    \end{equation}
%  }
%
%  By multiplying every term by 2, and again, moving $m_1$ terms to the left,
%  and $v_2$ terms to the right.
%\end{frame}



\begin{frame}{Conservation of Momentum \& Energy in Elastic Collisions}
  Dividing the equations (2) by (1) from the last slide, we get:

  \eq{-.1in}{
    \frac{(2)}{(1)}\quad\to\quad
    \frac{m_1(v_1^2-v_1'^2)}{m_1(v_1-v_1')}=
    \frac{m_2(v_2'^2-v_2^2)}{m_2(v_2'-v_2)}
  }

  $m_1$ and $m_2$ terms cancel out, while the terms in the numerator can be
  expanded as the difference of two squares which is then simplified:

  \eq{-.1in}{
    \frac{(v_1+v_1')(v_1-v_1')}{(v_1-v_1')}=
    \frac{(v_2'+v_2)(v_2'-v_2)}{(v_2'-v_2)}
  }
  
  Leading to the final expression, which is substituted back into (1)
  
  \eq{-.1in}{
    v_1 +v_1'= v_2+v_2'
  }

\end{frame}



%\begin{frame}{Solving the Example Problem}
%  
%  {\Large
%    \begin{displaymath}
%      \boxed{v_A'=\frac{m_A-m_B}{m_A+m_B}v_A}\quad
%      \boxed{v_B'=\frac{2m_A}{m_A+m_B}v_A}
%    \end{displaymath}
%  }
%
%  These equations work for \emph{all} elastic impact where object B (in this
%  example, the truck) is stationary when impact occurs. Substituting values for
%  $m_A$, $m_B$ and $v_A$, we get:
%  \begin{displaymath}
%    v_A'=\frac{m_A-m_B}{m_A+m_B}v_A=\frac{(1000-3000)}{(1000+3000)}\times 20
%    = \boxed{\SI{-10}{\metre\per\second}}
%  \end{displaymath}
%  \begin{displaymath}
%    v_B'=\frac{2m_A}{m_A+m_B}v_A=\frac{(2\times 1000)}{(1000+3000)}\times 20
%    = \boxed{\SI{10}{\metre\per\second}}
%  \end{displaymath}
%\end{frame}




\begin{frame}{Final Velocities in an Elastic Collision}
  When two objects 1 and 2 of mass $m_1$ and $m_2$ and  collide elastically,
  their final velocities will be determined by the masses ($m_1$, $m_2$) and
  initial velocities ($v_1$, $v_2$):
  
  \vspace{-.2in}{\large
    \begin{align*}
      v_1'&=\frac{m_1-m_2}{m_2+m_1}v_1+\frac{2m_2}{m_2+m_1}v_2\\
      v_2'&=\frac{2m_1}{m_2+m_1}v_1 + \frac{m_2-m_1}{m_2+m_1}v_2
    \end{align*}
  }

  These equations are \emph{not} provided in the AP exam equation sheet, which
  means that we are more interested in the behaviour qualitatively rather than
  quantitatively.
\end{frame}




\begin{frame}{Special Case: Equal Mass}
  If both objects have equal mass ($m_1=m_2=m$), then the equations simplifies
  to
  
  \vspace{-.2in}{\large
    \begin{align*}
      v_1'&=\frac{v_1(m-m)+2mv_2}{m+m}=v_2\\
      v_2'&=\frac{v_2(m-m)+2mv_1}{m+m}=v_1
    \end{align*}
  }
  
  The objects ``exchange'' their velocities during the collision.
\end{frame}



\begin{frame}{Special Case: Equal Mass}
  Additionally, if the second object is initially at rest ($v_2=0$), then:
  
  \vspace{-.2in}{\large
    \begin{align*}
      v_1'&=\frac{v_1(m-m)+2mv_2}{m+m}=0\\
      v_2'&=\frac{v_2(m-m)+2mv_1}{m+m}=v_1
    \end{align*}
  }

  All the momentum and energy from $m_1$ is transferred to $m_2$. Object 1
  stops all together, while object 2 continues with the initial momentum and
  velocity of Object 1.
\end{frame}



\begin{frame}{Special Case: Large Object Hitting Small Object}
  Another special case is when $m_1\gg m_2$ and $v_2=0$ (i.e.\ a large object
  colliding with a small stationary object) then we can effectively ``ignore''
  $m_2$:
  
  \vspace{-.2in}{\large
    \begin{align*}
      v_1'&=\frac{v_1(m_1-m_2)+2m_2v_2}{m_1+m_2}\approx
      \frac{m_1v_1}{m_1}=v_1\\
      v_2'&=\frac{v_2(m_2-m_1)+2m_1v_1}{m_1+m_2}\approx
      \frac{2m_1v_1}{m_1}=2v_1
    \end{align*}
  }

  Object 1 continues to move like nothing happened, but object 2 is pushed to
  move at \emph{twice} the initial speed of object 1.
\end{frame}



\begin{frame}{Special Case: Small Object Bouncing off Large Object}
  In the reverse case, if $m_1\ll m_2$, and $v_2=0$ (a small object colliding
  with a large stationary object), then we can ``ignore'' the $m_1$ term:
  
  \vspace{-.2in}{\large
    \begin{align*}
      v_1'&=\frac{v_1(m_1-m_2)+2m_2v_2}{m_1+m_2}\approx
      \frac{-m_2v_1}{m_2}=-v_1\\
      v_2'&=\frac{v_2(m_2-m_1)+2m_1v_1}{m_1+m_2}\approx 0
    \end{align*}
  }

  Object 1 bounces off object 2, and travels in the opposite direction with the
  same velocity magnitude, while object 2 does not move.
\end{frame}



\begin{frame}{Elastic Collision Example}
  \textbf{Example:} Blocks A and B have the same mass; A hits B with a speed
  of \SI{5.0}{\metre\per\second} while B is initially at rest. If the collision
  is elastic, what would be the final speed of these two objects?
\end{frame}



\begin{frame}{Elastic Collision Example}
  \textbf{Example:} Blocks A and B with the same mass; A has a velocity
  \SI{3.0}{\metre\per\second} to the east while B has
  \SI{2.0}{\metre\per\second} to the west. If the collision is elastic, after
  the collision, what would the velocity of the two blocks be?
\end{frame}



\begin{frame}{Elastic Collision Example}
  \textbf{Example:} Throw a ball to a really big wall, when the ball reaches
  the wall, it has a velocity \SI{10}{\metre\per\second} toward the wall. If
  the collision is elastic, what would the final velocity of the ball be?
\end{frame}



\begin{frame}{Elastic Collision Example}
  \textbf{Example:} Throw a ball with a velocity \SI{4.0}{\metre\per\second}
  toward a train with a velocity \SI{40}{\metre\per\second} toward the ball.
  If the collision is elastic, what would the final velocity of the ball be?
\end{frame}


\begin{frame}{Inelastic Collision: Calculating Energy Loss}
  \textbf{Example:} Two blocks A and B with mass \SI{2.0}{\kilo\gram}, block
  A hits B with velocity \SI{4.0}{\metre\per\second} while B is at rest.
  \begin{enumerate}[(a)]
  \item Suppose the collision is completely inelastic, what would the final
    velocity of A and B be?
  \item What is the loss of energy?
  \end{enumerate}
\end{frame}


%\begin{frame}{Elastic Collision Example}
%  \textbf{Example:} Blocks A and B have the same mass; A hits B with a speed
%  of \SI{5.0}{\metre\per\second} while B is initially at rest. If the collision
%  is elastic, what would be the final speed of these two objects?
%\end{frame}
%
%
%
%\begin{frame}{Elastic Collision Example}
%  \textbf{Example:} Blocks A and B with the same mass; A has a velocity
%  \SI{3.0}{\metre\per\second} to the east while B has
%  \SI{2.0}{\metre\per\second} to the west. If the collision is elastic, after
%  the collision, what would the velocity of the two blocks be?
%\end{frame}
%
%
%
%\begin{frame}{Elastic Collision Example}
%  \textbf{Example:} Throw a ball to a really big wall, when the ball reaches
%  the wall, it has a velocity \SI{10}{\metre\per\second} towards the wall. If
%  the collision is elastic, what would the final velocity of the ball be?
%\end{frame}
%
%
%
%\begin{frame}{Elastic Collision Example}
%  \textbf{Example:} Throw a ball with a velocity \SI{4.0}{\metre\per\second}
%  towards a train with a velocity \SI{40}{\metre\per\second} towards the ball.
%  If the collision is elastic, what would the final velocity of the ball be?
%\end{frame}
%
%
%\begin{frame}{Inelastic Collision: Calculating Energy Loss}
%  \textbf{Example:} Two blocks A and B with mass \SI{2.0}{\kilo\gram}, block
%  A hits B with velocity \SI{4.0}{\metre\per\second} while B is at rest.
%  \begin{enumerate}[(a)]
%  \item Suppose the collision is completely inelastic, what would the final
%    velocity of A and B be?
%  \item What is the loss of energy?
%  \end{enumerate}
%\end{frame}
%
>>>>>>> 4c148618b825181207cf76db81d003de72e8315b
\end{document}
